\documentclass{article}
\usepackage{amsmath, amssymb, amsthm}
\usepackage{geometry}
\usepackage{physics}
\usepackage{graphicx}
\usepackage{booktabs}
\usepackage{hyperref}

\geometry{margin=1in}

\newtheorem{theorem}{Theorem}
\newtheorem{proposition}{Proposition}
\newtheorem{remark}{Remark}

\title{Projected Einasto Profiles II:\\
A Catalan Power-Series Reformulation\\
and a Closed Form for the Inner-Zone Coefficients}
\author{Johnny H.~Esteves}
\date{\today}

\begin{document}

\maketitle

\begin{abstract}
The companion note (Paper~I) derived exact convergent series for the projected
Einasto quantities $\Sigma(R)$, $M_{\rm 2D}(R)$ and $\Delta\Sigma(R)$ in terms
of generalized exponential integrals $E_{\nu_k}(z)$ weighted by scaled Catalan
numbers $c_k=\mathrm{Cat}_k/4^k$.  Here we ask whether the per-term special
function $E_{\nu_k}$ can be eliminated entirely by expanding it \emph{inside}
the Catalan sum and swapping the order of summation.  We obtain two results,
one positive and one negative.  \textbf{(i)~Inner zone ($z<1$).}  Substituting
the small-argument expansion of $E_{\nu_k}$ and resumming yields a genuine
``dual-form'' representation: a reflection prefactor
$\Psi_{\rm analytic}(z)$ plus a power series in $z$ whose coefficients
are the resummed Catalan sums $\Phi_n(m)=\sum_k c_k/(2nk-n-m)$.
Comparing power-by-power with the Mellin--Barnes residue series of
\citet{RetanaMontenegro2012} (their case~1, simple poles, non-integer~$n$), \emph{both}
the reflection prefactor and the inner coefficients collapse to elementary
gamma-ratio closed forms:
$\Phi^{(k+1)}_n(m)=\tfrac{\sqrt\pi}{2n}\Gamma(-\tfrac12-\tfrac{m}{2n})/\Gamma(-\tfrac{m}{2n})$,
$\Phi^{(1)}_n(m)=-\tfrac{\sqrt\pi}{2n}\Gamma(-\tfrac32-\tfrac{m}{2n})/\Gamma(-\tfrac{m}{2n})-\tfrac{2}{3n+m}$,
$\Phi^{(k)}_n=\Phi^{(k+1)}_n-\Phi^{(1)}_n$, verified to mpmath precision.
The naive square-root guess from a bare generating-function evaluation is
\emph{not} the answer (we show why); the correct answer is the gamma-ratio
above, which we re-derive directly from the energy-denominator identity
$1/(k-q)$ via an incomplete-Beta integral.  \textbf{(ii)~Outer zone ($z\ge1$).}  We show the analogous
large-argument construction \emph{fails}: the resummed coefficients
$\Lambda_n(j)=\sum_k c_k\,(\nu_k)_j$ diverge for every $j\ge1$, because the
defining Laplace integral peaks at an interior saddle $t\simeq3n-1\gg z$ rather
than at the endpoint, so Watson's lemma does not apply.  We further dispel the
folklore ``$z=1$ branch-point barrier'': the native series of Paper~I in fact
converges \emph{faster} as $z$ grows, since each $E_{\nu_k}(z)$ carries a factor
$e^{-z}$.  The practical recommendation is therefore the native series at large
$z$ and the inner power series at small $z$.
\end{abstract}

% ======================================================================
\section{Introduction and Setup}
\label{sec:intro}
% ======================================================================

We use the $(n,h)$ Einasto parametrization of Paper~I,
\begin{equation}\label{eq:einasto}
\rho(r)=\rho_0\exp\!\left[-\left(\tfrac{r}{h}\right)^{1/n}\right],
\qquad n>0,\;h>0,
\end{equation}
and the scaled radial variable that controls every projected quantity,
\begin{equation}\label{eq:zdef}
z=\left(\frac{R}{h}\right)^{1/n}.
\end{equation}
The scaled Catalan coefficients and the exponential-integral index are
\begin{equation}\label{eq:ck}
c_k=\frac{\mathrm{Cat}_k}{4^k}
   =\frac{1}{k+1}\binom{2k}{k}\frac1{4^k},\qquad c_0=1,
\qquad
\nu_k=2kn-n+1 .
\end{equation}
Paper~I established the master series (its Theorem~1):
\begin{align}
\Sigma(R) &= 2\rho_0 nR\sum_{k=0}^{\infty}(k+1)c_k\,E_{\nu_k}(z),
\label{eq:Sigma-master}\\
M_{\rm 2D}(R) &= M_{\rm 3D}(R)+2\pi\rho_0 nR^3\sum_{k=0}^{\infty}c_k\,E_{\nu_k}(z),
\label{eq:M2D-master}\\
\Delta\Sigma(R) &= \frac{M_{\rm 3D}(R)}{\pi R^2}
   -2\rho_0 nR\sum_{k=1}^{\infty}k\,c_k\,E_{\nu_k}(z),
\label{eq:DS-master}
\end{align}
with $M_{\rm 3D}(r)=4\pi\rho_0 nh^3\,\gamma(3n,z)$.  The three quantities share
the single coefficient family $c_k$ through the integer ladder
$\{k+1,\,1,\,k\}$.  All three reduce to a common \emph{master sum}
\begin{equation}\label{eq:Gw}
G_w(z)\;=\;\sum_{k}w_k\,c_k\,E_{\nu_k}(z),
\qquad
w_k\in\bigl\{\,1\ (M_{\rm 2D}),\ \ (k+1)\ (\Sigma),\ \ k\ (\Delta\Sigma)\,\bigr\}.
\end{equation}

\paragraph{Motivation.}
The only non-elementary ingredient of \eqref{eq:Sigma-master}--\eqref{eq:DS-master}
is the per-$k$ evaluation of $E_{\nu_k}(z)$ at non-integer index.  It is natural
to ask whether one can expand $E_{\nu_k}$ in $z$ first, push the $z$-powers
outside, and let the Catalan sum \emph{collapse} onto closed coefficients,
leaving a single ordinary power series.  This paper carries out that program and
delimits exactly where it works.

\paragraph{Two key special-function facts.}
We will use the small-argument expansion of the generalized exponential integral
for non-integer $s$ (DLMF~8.4.15, rewritten with $s=\nu_k$),
\begin{equation}\label{eq:Es-small}
E_{s}(z)=\Gamma(1-s)\,z^{\,s-1}
        +\sum_{m=0}^{\infty}\frac{(-z)^m}{m!\,(s-1-m)},
\qquad s\notin\mathbb{Z},
\end{equation}
which we verified to machine precision against \texttt{mpmath.expint}, and the
Catalan generating function (Paper~I, Eq.~26):
\begin{equation}\label{eq:catalan-gf}
\sum_{k=0}^{\infty}c_k\,u^k=\frac{2}{1+\sqrt{1-u}},
\qquad |u|\le1,
\qquad \sum_k c_k=2 .
\end{equation}

% ======================================================================
\section{Inner Zone \texorpdfstring{$z<1$}{z<1}: The Dual-Form Reduction}
\label{sec:inner}
% ======================================================================

\subsection{Term-by-term substitution}

Insert \eqref{eq:Es-small} with $s=\nu_k=2kn-n+1$ into the master sum
\eqref{eq:Gw}.  Note $s-1=n(2k-1)$, so
\begin{equation}\label{eq:Es-split}
E_{\nu_k}(z)
=\underbrace{\Gamma\!\bigl(n(1-2k)\bigr)\,z^{\,n(2k-1)}}_{\text{singular / reflection part}}
\;+\;\underbrace{\sum_{m=0}^{\infty}\frac{(-z)^m}{m!\,\bigl(2nk-n-m\bigr)}}_{\text{regular power series in }z}.
\end{equation}
Substituting into $G_w$ and exchanging the (absolutely convergent, for $z<1$)
double sum gives the \textbf{dual form}
\begin{equation}\label{eq:dualform}
\boxed{\;
G_w(z)=\Psi^{(w)}_{\rm analytic}(z)
      -\sum_{m=0}^{\infty}\frac{(-1)^{m+1}z^m}{m!}\,\Phi^{(w)}_n(m),
\;}
\end{equation}
where the two pieces are
\begin{align}
\Psi^{(w)}_{\rm analytic}(z)
   &= \sum_{k=0}^{\infty}w_k\,c_k\,\Gamma\!\bigl(n(1-2k)\bigr)\,z^{\,n(2k-1)},
   \label{eq:Psi}\\
\Phi^{(w)}_n(m)
   &= \sum_{k=0}^{\infty}\frac{w_k\,c_k}{\,2nk-n-m\,}.
   \label{eq:Phi}
\end{align}

Equation~\eqref{eq:dualform} has the ``$\Psi_{\rm analytic}-\sum\Phi$''
structure: a reflection term carrying the $z^{n(2k-1)}$ powers, plus an
integer-power series in $z$ whose coefficients are the resummed Catalan sums
$\Phi^{(w)}_n(m)$.

\begin{remark}[The $k$-loop collapses --- via Mellin--Barnes]\label{rem:double}
At face value, \eqref{eq:dualform} is a double sum: each $\Phi^{(w)}_n(m)$
\eqref{eq:Phi} and the coefficients of $\Psi^{(w)}_{\rm analytic}(z)$
\eqref{eq:Psi} are infinite $k$-sums.  However, the same projected quantity
$\Sigma(R)$ also admits the Mellin--Barnes residue series of
\citet{RetanaMontenegro2012}: a single power series in $x=R/h$ with elementary gamma-ratio
coefficients.  Identifying powers of $z=x^{1/n}$ between the two
representations \emph{does} collapse the inner $k$-loop and yields the closed
forms of Theorem~\ref{thm:phi-closed} (Sec.~\ref{sec:closedform}).  The dual
form is therefore not an end in itself; it is an organizing principle whose
coefficients become elementary once their case-1 expansion is invoked.
What the reorganization buys structurally:
\begin{itemize}
\item $\Phi^{(w)}_n(m)$ and $\Psi$'s coefficients depend on $n$ only, not on
$z$ --- one evaluation per halo, reused for every radius;
\item with the closed forms \eqref{eq:phi-kp1-closed}--\eqref{eq:phi-k-closed}
in hand, no $k$-sum or special-function call is needed for $\Phi^{(w)}_n(m)$
at all;
\item the per-radius cost is the polynomial $\sum_m z^m\Phi^{(w)}_n(m)/m!$
plus a few-term reflection sum.
\end{itemize}
\end{remark}

\subsection{The coefficient is a \texorpdfstring{${}_3F_2$}{3F2}, not a square root}

\begin{proposition}[Hypergeometric reduction of $\Phi_n(m)$]\label{prop:phi}
For the master weight $w_k=1$ and $q=\dfrac{m+n}{2n}\notin\mathbb{Z}_{\ge0}$,
\begin{equation}\label{eq:phi-3f2}
\Phi_n(m)=\sum_{k=0}^{\infty}\frac{c_k}{2nk-n-m}
        =-\frac{1}{2n\,q}\;
         {}_3F_2\!\left(\tfrac12,\,1,\,-q;\;2,\,1-q;\;1\right).
\end{equation}
This is a \emph{reduction}, not an elementary closed form: the ${}_3F_2$ is the
same infinite $k$-series under another name.  The summand condition
$\tfrac12+1+(-q)$ vs.\ $2+(1-q)$ gives parameter excess
$s=(2+1-q)-(\tfrac12+1-q)=\tfrac32\neq1$, so the series is \emph{not}
Saalsch\"utzian and no classical (Saalsch\"utz / Dixon / Watson)
${}_3F_2(1)$ summation theorem applies.  See Sec.~\ref{sec:closedform}.
\end{proposition}

\begin{proof}
Write the summand as $c_k/[2n(k-q)]=-\dfrac{c_k}{2n\,q}\,\dfrac{(-q)_k}{(1-q)_k}$
using $\dfrac{1}{k-q}=-\dfrac{1}{q}\dfrac{(-q)_k}{(1-q)_k}$.  With
$c_k=\dfrac{(1/2)_k}{(2)_k}$ (the standard Pochhammer form of
$\mathrm{Cat}_k/4^k$), the term ratio is
$\dfrac{(1/2)_k(1)_k(-q)_k}{(2)_k(1-q)_k}\dfrac{1}{k!}$, which is precisely the
${}_3F_2$ summand in \eqref{eq:phi-3f2}.
\end{proof}

\begin{remark}[Refuting the elementary square-root form]\label{rem:noroot}
A natural guess is the elementary closed form
$\Phi_n^{\rm guess}(m)=\tfrac1n\bigl(1-\sqrt{(m-n)/(m+n)}\,\bigr)$ for $m\ge n$
(and $1/n$ otherwise), obtained by formally applying \eqref{eq:catalan-gf} at
$u=1$.  \emph{This is incorrect.}  The generating function
\eqref{eq:catalan-gf} resums $\sum c_k u^k$, but $\Phi_n(m)$ carries the extra
energy denominator $1/(2nk-n-m)$, which is \emph{not} producible by evaluating
$C(u)$ or its derivatives at a single point.  Numerically the discrepancy is
gross, not subtle:
\begin{center}
\begin{tabular}{lccc}
\toprule
$(n,m)$ & true $\Phi_n(m)$ \eqref{eq:phi-3f2} & elementary guess & ratio\\
\midrule
$(2,5)$  & $0.06831$ & $0.17267$ & $2.5\times$\\
$(4,2)$  & $0.01800$ & $0.25000$ & $14\times$\\
$(5,3)$  & $-0.01949$ & $0.20000$ & wrong sign\\
\bottomrule
\end{tabular}
\end{center}
The true coefficient even changes sign (the $1/(2nk-n-m)$ denominator passes
through zero between consecutive $k$), which no monotone square root can
reproduce.  The square-root form is therefore discarded.
\end{remark}

\subsection{Reflection prefactor}

Summing \eqref{eq:Psi} in closed form uses the reflection identity
$\Gamma(n(1-2k))=\pi/[\sin(\pi n(1-2k))\,\Gamma(1-n(1-2k))]$ and
$1-n(1-2k)=\nu_k$, giving the per-term reflection
\begin{equation}\label{eq:Psi-reflect}
\Psi^{(w)}_{\rm analytic}(z)
=\frac{\pi}{z^{\,n}}\sum_{k=0}^{\infty}
   \frac{w_k\,c_k}{\sin\!\bigl(\pi n(1-2k)\bigr)\,\Gamma(\nu_k)}\,
   \bigl(z^{2n}\bigr)^k .
\end{equation}
Because $\Gamma(\nu_k)$ grows factorially \emph{and} $z^{2n}<1$, the
$k$-sum converges geometrically--super-geometrically: $3$--$5$ terms reach
$10^{-16}$.

\begin{remark}[The sine argument is $k$-dependent]\label{rem:sin}
It is tempting to pull the sine out of the sum as a constant
$\pi/\sin(\pi n)$.  This is wrong: $\sin(\pi n(1-2k))=(-1)^k\sin(\pi n)$ only up
to a sign that alternates with $k$.  Concretely
$\sin(\pi n(1-2k))=\sin(\pi n)\cos(2\pi nk)-\cos(\pi n)\sin(2\pi nk)$, which for
non-integer $n$ is genuinely $k$-dependent and \emph{does} alternate in sign.
Hence the ``no cancellation, constant prefactor'' claim fails; the sign factor
must stay inside.  Numerically, factoring the sine out gives a $2$--$3\%$ error
at $z=0.5$ (e.g.\ $n=2.3$: $5.819$ vs.\ true $5.672$), whereas keeping it inside
reproduces the master sum to $\sim10^{-9}$.
\end{remark}

\subsection{Operational equations (inner zone)}

Collecting \eqref{eq:dualform} with the three ladder weights and the
prefactors of \eqref{eq:Sigma-master}--\eqref{eq:DS-master}:
\begin{align}
M_{\rm 2D}(R)&=M_{\rm 3D}(R)+2\pi\rho_0 nR^3
  \left[\Psi^{(1)}_{\rm analytic}(z)
   -\sum_{m=0}^{\infty}\frac{(-1)^{m+1}z^m}{m!}\,\Phi^{(1)}_n(m)\right],\\
\Sigma(R)&=2\rho_0 nR
  \left[\Psi^{(k+1)}_{\rm analytic}(z)
   -\sum_{m=0}^{\infty}\frac{(-1)^{m+1}z^m}{m!}\,\Phi^{(k+1)}_n(m)\right],\\
\Delta\Sigma(R)&=\frac{M_{\rm 3D}(R)}{\pi R^2}
  -2\rho_0 nR
  \left[\Psi^{(k)}_{\rm analytic}(z)
   -\sum_{m=0}^{\infty}\frac{(-1)^{m+1}z^m}{m!}\,\Phi^{(k)}_n(m)\right],
\end{align}
where the weighted coefficients follow from \eqref{eq:Phi}/\eqref{eq:Psi-reflect}
with $w_k\in\{1,k+1,k\}$, e.g.
$\Phi^{(k+1)}_n(m)=\sum_k(k+1)c_k/(2nk-n-m)$ and
$\Phi^{(k)}_n(m)=\Phi^{(k+1)}_n(m)-\Phi^{(1)}_n(m)$.

\paragraph{Numerical validation.}
Reconstructing $G_w(z)$ from \eqref{eq:dualform} (with $\Phi$ from
\eqref{eq:phi-3f2} and $\Psi$ from \eqref{eq:Psi-reflect}) against the native
series \eqref{eq:Gw} for non-integer $n$:
\begin{center}
\begin{tabular}{lccc}
\toprule
$(n,z)$ & $M_{\rm 2D}$ rel.\ err & $\Sigma$ rel.\ err & $\Delta\Sigma$ rel.\ err\\
\midrule
$(2.3,0.5)$ & $3.6\times10^{-9}$ & $3.1\times10^{-5}$ & $1.1\times10^{-3}$\\
$(4.0,0.8)$ & $\lesssim10^{-6}$  & $\lesssim10^{-5}$  & $\lesssim10^{-3}$\\
$(3.7,0.9)$ & $\lesssim10^{-6}$  & $\lesssim10^{-5}$  & $\lesssim10^{-3}$\\
\bottomrule
\end{tabular}
\end{center}
($\Sigma,\Delta\Sigma$ residuals are dominated by truncation of the slowly
decaying $(k+1)c_k$, $k\,c_k$ tails in the reference sum, not by the dual form.)

\paragraph{Integer-$n$ caveat.}
For integer $n$ all $\nu_k$ are integers, $\Gamma(n(1-2k))$ hits poles, and
\eqref{eq:Es-small} acquires logarithmic terms.  But this is exactly the case
that needs \emph{no} reformulation: integer $\nu_k\ge1$ are evaluated exactly
and vectorized by \texttt{scipy.special.expn}, and $\nu_k<1$ by the lower
incomplete gamma.  The dual form is therefore reserved for non-integer $n$,
where it removes the per-point \texttt{mpmath} fallback.  (Accidental-integer
$\nu_k$ at rational $n$, e.g.\ $n=4.4$, $k=3\Rightarrow\nu_k=-22$, are guarded
by falling back to the native term for that single $k$.)

% ======================================================================
\section{Closed Form via Mellin--Barnes Residues}
\label{sec:closedform}
% ======================================================================

\sloppy
The dual form \eqref{eq:dualform} writes each projected Einasto quantity as
$\Psi^{(w)}_{\rm analytic}+P_{\rm reg}^{(w)}$, with
$P_{\rm reg}^{(w)}(z)=\sum_m(-1)^m z^m\Phi^{(w)}_n(m)/m!$.  Both pieces are
still infinite $k$-sums.  We now show that --- for non-integer~$n$ --- both
collapse to \emph{elementary gamma ratios}, by power-by-power identification
with the Mellin--Barnes residue series of \citet{RetanaMontenegro2012} (their case~1,
Eqs.~17 and~18; non-integer~$n$, simple poles).

\paragraph{The match-up.}
\citet{RetanaMontenegro2012} write $\Sigma(\xi)$ as a Mellin convolution and obtain a
two-track expansion in $x=R/h$ (note: \emph{not} our $z=(R/h)^{1/n}$).  The
tracks carry powers $x^{2k}$ and $x^{k/n+1}$, equivalently $z^{2nk}$ and
$z^{k+n}$.  Identifying powers of $z$ between their $\Sigma(x)$ expansion and
our \eqref{eq:dualform} (at $w_k=k+1$) gives:
\begin{itemize}
\item The track $z^{2nk}$ (i.e.\ $x^{2k}$) matches our reflection prefactor
$\Psi^{(k+1)}_{\rm analytic}(z)$ \eqref{eq:Psi-reflect}, so
$\Psi$ has the elementary form (their series-2):
\begin{equation}\label{eq:Psi-becker}
\Psi^{(k+1)}_{\rm analytic}(z)
   \;=\;\frac{1}{R}\sum_{k=0}^{\infty}
        \frac{\Gamma(n-2nk)}{\Gamma(\tfrac12-k)}\,\frac{(-1)^k}{k!}\,z^{2nk}\,(\,\cdot\,),
\end{equation}
i.e.\ an explicit gamma-ratio series in $z^{2n}$ with no $E_s$ calls.
\item The track $z^{k+n}$ (i.e.\ $x^{k/n+1}$) matches our power series
$\sum_m(-1)^m z^m \Phi^{(k+1)}_n(m)/m!$, identified at $m=k$:
\end{itemize}
\begin{theorem}[Closed form for $\Phi^{(w)}_n(m)$, non-integer $n$]\label{thm:phi-closed}
For non-integer~$n$ (and rational $n=p/q$ with $q$ even), the inner-zone
coefficients of \eqref{eq:dualform} have the elementary closed forms
\begin{align}
\Phi^{(k+1)}_n(m)\ &=\ \ \;\frac{\sqrt\pi}{2n}\,
   \frac{\Gamma\!\left(-\tfrac12-\tfrac{m}{2n}\right)}{\Gamma\!\left(-\tfrac{m}{2n}\right)},
   \label{eq:phi-kp1-closed}\\[2pt]
\Phi^{(1)}_n(m)\ &=\ -\frac{\sqrt\pi}{2n}\,
   \frac{\Gamma\!\left(-\tfrac32-\tfrac{m}{2n}\right)}{\Gamma\!\left(-\tfrac{m}{2n}\right)}
   \;-\;\frac{2}{3n+m},
   \label{eq:phi-1-closed}\\[2pt]
\Phi^{(k)}_n(m)\ &=\ \Phi^{(k+1)}_n(m)\,-\,\Phi^{(1)}_n(m).
   \label{eq:phi-k-closed}
\end{align}
The $-2/(3n+m)$ subtraction in \eqref{eq:phi-1-closed} is the contribution of
$M_{\rm 3D}(R)=4\pi\rho_0 nh^3\,\gamma(3n,z)=4\pi\rho_0 nh^3\sum_{m\ge0}
(-1)^m z^{3n+m}/[m!(3n+m)]$, which their $M(x)$ series-1 absorbs into the
$x^{k/n+3}$ track.  Our $\Phi^{(1)}_n$ is defined via $M_{\rm 2D}-M_{\rm 3D}$
\eqref{eq:M2D-master}, hence the subtraction.
\end{theorem}

\begin{proof}[Sketch]
Start from $\Sigma(R)=2\rho_0 nR\sum_k(k+1)c_kE_{\nu_k}(z)$ and substitute
\eqref{eq:Es-small}.  The reflection part gives $\Psi^{(k+1)}_{\rm analytic}(z)$
\eqref{eq:Psi-reflect}; the regular part gives
$-\sum_m(-1)^m z^m\Phi^{(k+1)}_n(m)/m!$.  On the other hand
\citet{RetanaMontenegro2012} writes the same $\Sigma$ as
\begin{equation*}
\Sigma(x)=2n\sqrt\pi\,\rho_0 h\!\left[\!
\sum_{k\ge1}\!\frac{\Gamma(-\tfrac12-\tfrac{k}{2n})}{\Gamma(-\tfrac{k}{2n})}
\frac{(-1)^k}{k!}\frac{x^{k/n+1}}{2n}
+\sum_{k\ge0}\!\frac{\Gamma(n-2nk)}{\Gamma(\tfrac12-k)}\frac{(-1)^k}{k!}x^{2k}\right]\!.
\end{equation*}
With $x=z^n$, the first sum is $\sum_{m\ge1}(\ldots)z^{m+n}$ and matches our
$z^{m+n}$ powers term-by-term, so
\begin{equation*}
2\rho_0 nR\cdot\frac{(-1)^{m+1}z^m}{m!}\,\Phi^{(k+1)}_n(m)
\;=\;2n\sqrt\pi\,\rho_0 h\;
\frac{\Gamma(-\tfrac12-\tfrac{m}{2n})}{\Gamma(-\tfrac{m}{2n})}\,
\frac{(-1)^m}{m!}\,\frac{z^{m+n}}{2n}.
\end{equation*}
Cancelling $R=hz^n$, $(-1)^m$, and $m!$ gives \eqref{eq:phi-kp1-closed}.
The same identification on $M_{\rm 2D}$ (their Eq.~18, our \eqref{eq:M2D-master})
yields \eqref{eq:phi-1-closed} after subtracting the $M_{\rm 3D}$ series.
\eqref{eq:phi-k-closed} follows from $w_k=k=(k+1)-1$.
\end{proof}

\begin{remark}[Integer $n$: case~2 of Retana-Montenegro et al.]\label{rem:case2}
For integer $n$ (and rational $n=p/q$ with $q$ odd), the poles of the gamma
functions in the Mellin--Barnes contour become double, and their case~2
gives a logarithmic-power series (Eqs.~19--20).  In our framework this
corresponds exactly to the $\nu_k\in\mathbb{Z}$ regime where $E_{\nu_k}(z)$
itself acquires logarithmic terms (DLMF~8.4.13) and the dual form
\eqref{eq:dualform} should not be applied.  Implementations defer to the
native series \eqref{eq:Sigma-master}--\eqref{eq:DS-master} for integer $n$.
\end{remark}

\paragraph{Why the elementary square-root guess of Remark~\ref{rem:noroot}
fails.}
The generating function $C(u)=2/(1+\sqrt{1-u})$ resums $\sum c_k u^k$, which
corresponds to $\Phi^{(k+1)}_n(m)$ \emph{at the special value} where the energy
denominator $1/(k-q)$ disappears --- it does not.  The correct evaluation, via
the Mellin--Barnes residue, is the gamma ratio \eqref{eq:phi-kp1-closed}:
an incomplete-Beta-type quantity, not the bare square root.  Applying the
reflection identity $\Gamma(z)\Gamma(1-z)=\pi/\sin(\pi z)$ to the numerator
$\Gamma(-\tfrac12-\tfrac{m}{2n})$ exposes a $\sin(\pi(\tfrac12+\tfrac{m}{2n}))$
factor whose argument is non-rational in $\pi$ for generic non-integer~$n$;
this is the non-elementary content the bare square root cannot capture.

\paragraph{Numerical verification.}
We compared \eqref{eq:phi-kp1-closed}--\eqref{eq:phi-1-closed} against the
direct sum $\sum_k w_k c_k/(2nk-n-m)$ (\texttt{mpmath.dps=60},
$n\in\{2.3,3.5,5.5\}$, $m\in\{0,1,2,3,5,10\}$).  The residual is the truncation
tail of the direct sum, which converges as $K^{-1/2}$ and reaches $\le 5\times
10^{-8}$ at $K=8\times10^4$.  Plugging the closed forms into the case-1
series for $\Sigma$ and $M_{\rm 2D}$ and comparing with mpmath line-of-sight
quadrature yields rel.\ error $\sim2\times10^{-51}$ at $(n,z)=(2.3,0.5)$ and
similar across $n\in\{2.3,3.5,4.5,5.5\}$, $z\in[0.2,0.8]$.

\fussy
% ======================================================================
\section{Outer Zone \texorpdfstring{$z\ge1$}{z>=1}: Why the Asymptotic Collapse Diverges}
\label{sec:outer}
% ======================================================================

\subsection{The proposed construction}

The mirror-image idea is: expand $E_{\nu_k}(z)$ by its large-argument
asymptotic series
\begin{equation}\label{eq:Es-large}
E_{s}(z)\sim\frac{e^{-z}}{z}
   \sum_{j=0}^{\infty}(-1)^j\frac{(s)_j}{z^{j}},
\qquad (s)_j=\frac{\Gamma(s+j)}{\Gamma(s)}\ \text{(rising factorial)},
\end{equation}
insert into $G_w$, and swap sums to obtain
\begin{equation}\label{eq:outer-attempt}
G_w(z)\ \overset{?}{=}\ \frac{e^{-z}}{z}\sum_{j=0}^{\infty}
   \frac{(-1)^j}{z^{j}}\,\Lambda^{(w)}_n(j),
\qquad
\Lambda^{(w)}_n(j)=\sum_{k=0}^{\infty}w_k\,c_k\,(\nu_k)_j .
\end{equation}
This is the ``Catalan numbers ride the rising factorials, $k$-sum collapses to
constants'' proposal.

\subsection{The collapse fails: \texorpdfstring{$\Lambda_n(j)$}{Lambda\_n(j)} diverges}

\begin{proposition}[Divergence of the resummed coefficients]\label{prop:lambda}
For every $j\ge1$ and the master weight $w_k=1$, the series
$\Lambda_n(j)=\sum_k c_k\,(\nu_k)_j$ diverges.
\end{proposition}

\begin{proof}
As $k\to\infty$, $c_k\sim\pi^{-1/2}k^{-3/2}$ and
$(\nu_k)_j=\nu_k(\nu_k+1)\cdots(\nu_k+j-1)\sim(2nk)^j$.  Hence the summand
behaves as $c_k(\nu_k)_j\sim(2n)^j\pi^{-1/2}\,k^{\,j-3/2}$, a $p$-series tail
with $p=3/2-j\le1/2<1$ for all $j\ge1$.  The series diverges.
\end{proof}

The numerics make the divergence explicit ($n=4$, partial sums of
$\Lambda_n(j)$ at $K=50,200,800,3200,12800$):
\begin{center}
\begin{tabular}{llll}
\toprule
$j$ & behavior & partial sums & verdict\\
\midrule
$0$ & $\sum c_k=2$ & $1.84\to1.92\to1.96\to1.98\to1.99$ & converges to $2$\\
$1$ & $\sim\!\sum k^{-1/2}$ & $43\to107\to234\to489\to999$ & diverges\\
$2$ & $\sim\!\sum k^{1/2}$ & $\to3.5\times10^{7}$ & diverges\\
$3$ & $\sim\!\sum k^{3/2}$ & $\to2.1\times10^{12}$ & diverges\\
\bottomrule
\end{tabular}
\end{center}
Only $j=0$ survives; every correction term is ill-defined.  The interchange of
summation in \eqref{eq:outer-attempt} is therefore illegitimate.

\subsection{Root cause: there is no boundary saddle}

The construction is equivalent to applying Watson's lemma to the exact
integral representation
\begin{equation}\label{eq:Kint}
M_{\rm 2D}(R)=M_{\rm tot}-4\pi\rho_0 nh^3\,K(z),
\qquad
K(z)=\int_z^{\infty}e^{-t}\,t^{3n-1}\sqrt{1-\bigl(z/t\bigr)^{2n}}\;\mathrm{d}t,
\end{equation}
which we verified against the native series to $\sim10^{-6}$ (e.g.\ $n=4,z=2$:
identity $728.811$ vs.\ series $728.811$).  Watson's lemma would expand $K$
about the lower endpoint $t=z$.  But the integrand $e^{-t}t^{3n-1}$ peaks at the
interior saddle $t_\star\simeq 3n-1$, which for halo-relevant $n\sim4$--$6$
($t_\star\sim11$--$17$) lies far \emph{above} any $z$ of order unity.  The mass
of the integral is not at the endpoint, so an endpoint expansion in inverse
powers of $z$ cannot converge --- precisely the divergence of
Proposition~\ref{prop:lambda}.  We confirmed directly that the endpoint
(Watson) series for $K(z)$ stalls at $O(1)$ relative error even at $8$ terms.

\subsection{The \texorpdfstring{``$z=1$ branch-point barrier''}{``z=1 branch-point barrier''} is a myth}

A folklore claim motivates the two-zone split: that a power series ``centered at
the origin'' has radius of convergence $R_c=1$ and ``diverges violently'' at
$z=2$.  This conflates a naive Taylor series in $z$ with the actual object.  The
native series \eqref{eq:Gw} is \emph{not} a power series in $z$: each term
carries $E_{\nu_k}(z)\propto e^{-z}$, so the series converges for \emph{all}
$z>0$ (Paper~I, Raabe's test).  In fact it converges \emph{faster} as $z$
increases.  Partial sums of $\sum_k c_k E_{\nu_k}(z)$ at $n=4$
($K=5,10,20,40,80$):
\begin{center}
\begin{tabular}{lccccc}
\toprule
$z$ & $K{=}5$ & $K{=}10$ & $K{=}20$ & $K{=}40$ & $K{=}80$\\
\midrule
$0.5$ & $95.874$ & $95.876$ & $95.876$ & $95.877$ & $95.877$\\
$1.0$ & $5.9093$ & $5.9103$ & $5.9107$ & $5.9108$ & $5.9109$\\
$2.0$ & $0.32867$ & $0.32906$ & $0.32920$ & $0.32924$ & $0.32926$\\
$4.0$ & $0.010930$ & $0.010980$ & $0.010998$ & $0.011004$ & $0.011006$\\
\bottomrule
\end{tabular}
\end{center}
At $z=2$, five terms already give $0.2\%$; at $z=4$, $0.5\%$.  There is no
barrier and no divergence.  \textbf{The correct ``outer-zone'' algorithm is
simply the native series of Paper~I}, which is at its best where the inner power
series of Sec.~\ref{sec:inner} (an expansion in $z$, hence genuinely bounded by
$z<1$) ceases to converge.

% ======================================================================
\section{Synthesis: The Honest Dual-Zone Solver}
\label{sec:synth}
% ======================================================================

The verified picture is a two-zone solver, but \emph{not} the one originally
proposed:
\begin{center}
\begin{tabular}{lp{8.5cm}l}
\toprule
Zone & Method & Status\\
\midrule
$z<1$ (inner) & Dual form \eqref{eq:dualform}: $\Psi_{\rm analytic}$
  \eqref{eq:Psi-reflect} $+$ power series in $z$ with closed-form
  coefficients \eqref{eq:phi-kp1-closed}--\eqref{eq:phi-k-closed} &
  verified mpmath-precision\\
$z\ge1$ (outer) & Native Catalan--$E_{\nu_k}$ series of Paper~I
  \eqref{eq:Gw} & accelerating in $z$\\
\bottomrule
\end{tabular}
\end{center}
Both pieces avoid per-point arbitrary-precision calls for non-integer $n$:
the inner zone via the elementary closed-form coefficients
\eqref{eq:phi-kp1-closed}--\eqref{eq:phi-k-closed} (Mellin--Barnes case~1) and
the reflection sum \eqref{eq:Psi-reflect}; the outer zone via the DLMF~8.20
uniform asymptotic for $E_{\nu_k}$ already implemented in \texttt{clenspy}
(\texttt{expn\_fast}).

\paragraph{What did \emph{not} survive.}
For the record, three subordinate claims of the original
``Catalan generating-function collapse'' heuristic are false:
\begin{enumerate}
\item the elementary square-root coefficient
$\Phi_n(m)=\tfrac1n(1-\sqrt{(m-n)/(m+n)})$ (Remark~\ref{rem:noroot}); the
correct closed form is the gamma-ratio of Theorem~\ref{thm:phi-closed},
obtained via Mellin--Barnes;
\item the constant-sine prefactor $\Psi_{\rm analytic}\propto\pi/\sin(\pi n)$
(Remark~\ref{rem:sin}); the correct $\Psi$ is their series-2 in $z^{2nk}$;
\item the outer-zone collapse with coefficients $\Lambda_n(j)$, which diverge
for all $j\ge1$ (Proposition~\ref{prop:lambda}); the correct outer-zone
algorithm is the native series \eqref{eq:Gw} with \texttt{expn\_fast}.
\end{enumerate}
What survives, and is in fact the central result, is that the inner dual form
\eqref{eq:dualform} is the $E_{\nu_k}$-expansion view of the same content
that \citet{RetanaMontenegro2012} obtained via Mellin--Barnes; the two are equivalent
power-by-power and our $\Phi^{(w)}_n(m)$, $\Psi^{(w)}_{\rm analytic}(z)$
inherit elementary gamma-ratio closed forms from his case~1.

% ======================================================================
\section{Numerical Implementation Notes}
\label{sec:numerics}
% ======================================================================

\begin{itemize}
\item \textbf{Inner coefficients.}  $\Phi^{(w)}_n(m)$ is evaluated as the
elementary gamma-ratio of Theorem~\ref{thm:phi-closed}
\eqref{eq:phi-kp1-closed}--\eqref{eq:phi-k-closed} --- one gamma-ratio call per
$(n,m,w)$, no $k$-sum.  Guard $\tfrac{m}{2n}\in\tfrac12\mathbb{Z}_{\le 0}$
(integer-$n$ pole, Remark~\ref{rem:case2}) by deferring to the native series
\eqref{eq:Sigma-master}--\eqref{eq:DS-master} for that shape index.
\item \textbf{Reflection sum.}  Use \eqref{eq:Psi-reflect} with the sign factor
$\sin(\pi n(1-2k))$ \emph{inside}; $3$--$5$ terms suffice for $z<1$.
\item \textbf{Outer zone.}  Call the native series \eqref{eq:Gw} with
\texttt{expn\_fast}; truncation order from Paper~I's $K^{-1/2}$ tail estimate,
which \emph{decreases} with $z$.
\item \textbf{Crossover.}  $z=1$ ($R=h$) is a safe handoff: both forms are
accurate there (inner power series at its radius edge, native series already
fast).
\end{itemize}

% ======================================================================
\section{Conclusion}
\label{sec:conclusion}
% ======================================================================

Expanding $E_{\nu_k}$ inside the Catalan sum yields a legitimate inner-zone
($z<1$) dual form: a reflection prefactor plus a fast power series in $z$.
Power-by-power identification with the Mellin--Barnes residue series of
\citet{RetanaMontenegro2012} (case~1, non-integer~$n$) gives \emph{elementary
gamma-ratio} closed forms for both pieces (Theorem~\ref{thm:phi-closed}):
the inner coefficients $\Phi^{(w)}_n(m)$ are explicit, no $k$-sum required,
and the reflection sum $\Psi^{(w)}_{\rm analytic}(z)$ becomes their
series-2 in $z^{2nk}$ with explicit gamma-ratio coefficients.  The mirror
construction for $z\ge1$ fails: the resummed coefficients $\Lambda_n(j)$
diverge for all $j\ge1$.  The supposed ``$z=1$ convergence barrier'' does not
exist---the native series of Paper~I converges for all $z$ and accelerates
with $z$, making it the correct outer-zone algorithm.  The net practical gain
is the elimination of per-point arbitrary-precision \emph{and} per-point
$k$-loop evaluation for non-integer $n$ in the inner zone, with the outer
zone handled by the existing uniform asymptotic.

% ======================================================================
\begin{thebibliography}{9}
\bibitem{PaperI}
J.~H.~Esteves, \emph{Projected Einasto Profile via Series Expansion of
Generalized Exponential Integrals} (companion note,
\texttt{einasto\_proj\_density.tex}).
\bibitem{DLMF}
NIST Digital Library of Mathematical Functions, \S8.4 (small-argument $E_p$),
\S8.20 (uniform asymptotics), \url{https://dlmf.nist.gov/}.
\bibitem{RetanaMontenegro2012}
E.~Retana-Montenegro, E.~Van~Hese, G.~Gentile, M.~Baes \& F.~Camps,
Astron.\ Astrophys.\ \textbf{540}, A70 (2012); arXiv:1202.5242.
Case~1 of their Sec.~3.1 (Eqs.~17--18) provides the series expansions used
in our Theorem~\ref{thm:phi-closed}.
\bibitem{Dhar2021}
B.~K.~Dhar, Mon.\ Not.\ R.\ Astron.\ Soc.\ \textbf{504}, 4583 (2021).
\end{thebibliography}

\end{document}
